\Needspace{9\baselineskip}
\originalchapter{偏振、偏振片与波片}\label{ch:polar}
\originalsection{电场端点描出的曲线}
沿 $z$ 传播的电磁波, 电场只能在横向 $xy$ 平面变化. 固定空间一点观察电场端点, 可能沿一条直线往返, 也可能绕椭圆旋转. 这就是偏振的几何含义. 采用负时间指数, 复横振幅组成 Jones 向量 $\boldsymbol J=(E_{x0},E_{y0})^{\mathsf T}$, 实电场为 $\Rea[\boldsymbol J\e^{\ii(kz-\omega t)}]$. 整体相位改变时间起点, 不改变偏振椭圆.

取 $E_x=a\cos\chi$, $E_y=b\cos(\chi+\delta)$, $a,b\ge0$, $\chi=kz-\omega t$. 当 $a,b>0$ 且 $\sin\delta\ne0$, 消去 $\chi$ 得
\[
\frac{E_x^2}{a^2}+\frac{E_y^2}{b^2}-2\frac{E_xE_y}{ab}\cos\delta=\sin^2\delta.
\]
这是一个椭圆. $\delta=0$ 或 $\pi$ 时退化成直线, 另有一个分量为零也为线偏振. 当 $a=b$, $\delta=\pm\pi/2$ 时是圆偏振. 两分量相位差非零不等于一定圆偏振, 两分量等幅也不等于一定线偏振.

圆偏振的左右手名称在不同领域会采用不同视向约定. 本册直接用向量给出方向: $\boldsymbol J=E_0(1,\ii)^{\mathsf T}$ 在 $z=0$ 时 $\boldsymbol E=E_0(\cos\omega t,\sin\omega t)$, 随时间从 $+x$ 转向 $+y$; $(1,-\ii)$ 则从 $+x$ 转向 $-y$. 使用名称时必须同时说明观察方向, 避免仅凭“顺时针”判断. 平均强度与 $\boldsymbol J^\dagger\boldsymbol J=a^2+b^2$ 成正比.

\originalsection{偏振片的投影作用}
理想线偏振片沿易透射方向 $\boldsymbol u_\theta=(\cos\theta,\sin\theta)^{\mathsf T}$ 保留分量, 消除垂直分量. Jones 矩阵就是正交投影
\[
P_\theta=\boldsymbol u_\theta\boldsymbol u_\theta^{\mathsf T}
=\begin{pmatrix}\cos^2\theta&\cos\theta\sin\theta\\\cos\theta\sin\theta&\sin^2\theta\end{pmatrix},
\qquad \boldsymbol J_{\rm out}=P_\theta\boldsymbol J_{\rm in}.
\]
线偏振入射光与透射轴夹角 $\alpha$ 时, 输出电场幅度为输入的 $\cos\alpha$, 因而 $I_{\rm out}=I_{\rm in}\cos^2\alpha$, 称为马吕斯定律. 金属丝栅允许电子沿丝移动, 平行丝的电场容易激发电流而被反射或吸收, 所以透射易轴垂直于金属丝.

非偏振光是许多不相关发射过程的统计混合, 不能由一个固定 Jones 向量表示. 它并不是 $\boldsymbol J=(0,0)$, 后者根本没有光. 若横向平均强度各向同性, 投影到任一线轴的平均强度为原来的一半. 两个正交理想偏振片阻断光; 中间再插入角度 $\theta$ 的第三片, 原已通过第一片的线偏振光可有 $I_{\rm out}=I_1\cos^2\theta\sin^2\theta$. 投影改变了后续光的状态, 不是简单恢复原来被吸收的分量.

\originalsection{双折射与相位延迟}
波片的两条主轴具有不同折射率 $n_f,n_s$, 其中 $n_f<n_s$ 分别为快轴和慢轴. 厚度 $d$ 内的相位延迟为 $\Delta=\omega(n_s-n_f)d/c$. 除去快轴整体传播相位后, 对坐标轴重合的波片有 $W_0=\operatorname{diag}(1,\e^{\ii\Delta})$. 波片理想无损时保持总强度, 却改变相对相位.

快轴与 $x$ 轴夹角 $\theta$ 时, 先投影到主轴再传播, 最后转回实验室坐标:
\[
W_\theta=R_\theta W_0R_\theta^{\mathsf T},\quad
R_\theta=\begin{pmatrix}\cos\theta&-\sin\theta\\\sin\theta&\cos\theta\end{pmatrix}.
\]
四分之一波片取 $\Delta=\pi/2$（模 $2\pi$）, 使两等幅同相分量变为圆偏振, 但入射线偏振必须与快轴成 $45^\circ$ 才给等幅. 具体矩阵为
\[
W_\theta=\begin{pmatrix}
\cos^2\theta+\ii\sin^2\theta&(1-\ii)\sin\theta\cos\theta\\
(1-\ii)\sin\theta\cos\theta&\sin^2\theta+\ii\cos^2\theta
\end{pmatrix}.
\]
半波片取 $\Delta=\pi$, 将入射线偏振方向角 $\alpha$ 变成 $2\theta-\alpha$, 可由矩阵直接作用于 $(\cos\alpha,\sin\alpha)$ 验证. 波片厚度和折射率差决定设计波长; 对宽频光, 相位延迟并非处处相同. 多元件的矩阵按光通过的先后从右向左相乘, 顺序不能任意交换.

\begin{example}
$\boldsymbol J=E_0(1,1)^{\mathsf T}/\sqrt2$ 通过快轴沿 $x$ 的四分之一波片后, 再通过角度 $\theta$ 的线偏振片. 求两次输出的偏振和强度.
\end{example}
\begin{solution}
波片后为 $E_0(1,\ii)^{\mathsf T}/\sqrt2$, 强度不变, 电场圆偏振. 投影振幅 $\boldsymbol u_\theta^{\mathsf T}\boldsymbol J=E_0(\cos\theta+\ii\sin\theta)/\sqrt2$, 模平方为 $E_0^2/2$, 所以偏振片输出沿 $\boldsymbol u_\theta$ 线偏振, 强度为原来的 $1/2$, 与 $\theta$ 无关. 圆偏振每个时间的电场方向都在旋转, 对任一固定轴的平均投影相同.
\end{solution}
